\chapter{Conversion, Engines, and the Virtual Token}
\label{ch:vtok}
% ===================================================================

Chapter~\ref{ch:carho} left a question open and then made it unavoidable.
Once phlogiston is decomposed into a spectrum, a computation holds
balances in several kinds at once, and the natural question is whether
those kinds have a common measure --- whether, having so much storage and
so much bandwidth and so much authority, one can say how much the
computation \emph{has}.

The answer is that a common measure exists exactly when the ways of
converting one kind into another contain no free lunch, and that when it
exists it is essentially unique.
This chapter states that, because the ecology chapters need it: a
scientist eats, and eating converts what was inside another computation
into something this one can spend.
Part~\ref{part:physics} states it again, at length, and reads the same
theorem as a statement about energy.
Here it is only a statement about exchange rates.

\section{Conversion systems}

Work over the signature monoid $(\Sigma, \ast, ())$ of
Chapter~\ref{ch:carho}, whose elements index the layers of a token stack.

\begin{definition}[Conversion system]
\label{def:conv-system}
A \emph{conversion system} over $\Sigma$ is a triple
$\hypK = (\Tok, E, r)$ where $\Tok \subseteq \Sigma$ is a finite set of
token types; $E \subseteq \Tok \times \Tok$ is a symmetric set of
\emph{engines}, presented with both orientations; and
$r : E \to \Rpos$ is a \emph{rate} satisfying
\[
  r(a \to b) \cdot r(b \to a) = 1 .
\]
An engine $a \to b$ is read operationally: one unit of the $a$-token is
consumed and $r(a \to b)$ units of the $b$-token released.
We write $\Gamma_{\hypK} = (\Tok, E)$ for the underlying graph, which we
may take connected without loss of generality.
\end{definition}

The paradigm engine is a persistent exchange process --- a contract
sitting on a channel that swaps one kind of token for another and back
again.
That such a thing is expressible as an ordinary term, rather than as a
privileged operation of the runtime, is the whole point of stacks being
terms.
An engine is a computation, and it can be eaten like any other.

\begin{definition}[Monoidal compatibility]
\label{def:monoidal-conv}
$\hypK$ is \emph{monoidal} when the rates respect composition of
signatures: whenever $a \ast a'$ and $b \ast b'$ are tokens and the
component engines exist,
\[
  r(a \ast a' \to b \ast b') = r(a \to b) \cdot r(a' \to b') .
\]
\end{definition}

Monoidal compatibility says that the worth of a compound authority is
built from the worths of its components in the same way that $\ast$
builds the authority.

\section{No arbitrage is the virtual token}

Pass to logarithms.
The \emph{log-rate} is the edge $1$-cochain
$\ell(a \to b) := \log r(a \to b)$, antisymmetric by reversibility.

\begin{definition}[Arbitrage]
\label{def:arbitrage}
A \emph{cycle} is a closed edge path
$a_0 \to a_1 \to \cdots \to a_k = a_0$; its \emph{yield} is the product
of the rates along it.
The system exhibits \emph{arbitrage} if some cycle has yield $\neq 1$,
and is \emph{arbitrage-free} if every cycle has yield $1$.
\end{definition}

\begin{definition}[Virtual token]
\label{def:virtual-token}
A \emph{valuation}, or \emph{virtual token}, for $\hypK$ is a function
$\nu : \Tok \to \Rpos$ with
\[
  r(a \to b) = \frac{\nu(a)}{\nu(b)} \qquad \text{for every engine } a \to b .
\]
Writing $p = \log \nu$, the condition is $\ell = -\,\delta p$: the
log-rate is minus the graph gradient of a potential.
\end{definition}

\begin{theorem}[The fundamental equivalence]
\label{thm:vtok-ftap}
For a connected conversion system $\hypK$ the following are equivalent:
\begin{enumerate}[label=\textup{(\roman*)}]
  \item $\hypK$ is arbitrage-free;
  \item $\ell$ is exact, $\ell = -\,\delta p$ for some $p : \Tok \to \Real$;
  \item a valuation $\nu = \exp p$ exists.
\end{enumerate}
The potential, and hence the valuation, is unique up to an additive
(respectively multiplicative) constant --- a global choice of unit.
If $\hypK$ is monoidal, $\nu$ may be taken to be a monoid homomorphism,
and is then unique up to a character of $\Sigma$.
\end{theorem}

\begin{proof}
(i)$\Rightarrow$(ii): fix a basepoint $a_\star$ with $p(a_\star) = 0$ and
define $p(a)$ by summing $-\ell$ along any path from $a_\star$ to $a$.
Path-independence is exactly the vanishing of $\ell$ around cycles, so
$p$ is well defined, and $\ell = -\delta p$ by construction.
(ii)$\Rightarrow$(iii) is $\nu = \exp p$.
(iii)$\Rightarrow$(i) telescopes: a cycle has yield
$\prod_i \nu(a_i)/\nu(a_{i+1}) = 1$.
Uniqueness up to a constant is connectedness of $\Gamma_{\hypK}$.
For the monoidal refinement, additivity of $p$ on generators propagates
along $\ast$-compatible engines, leaving as residual freedom a
homomorphism $\Tok \to \Real$.
\end{proof}

\begin{remark}[The direction of explanation]
\label{rmk:vtok-direction}
Theorem~\ref{thm:vtok-ftap} is the combinatorial core of the fundamental
theorem of asset pricing, and what matters here is which way the
implication is read.
The virtual token does not have to be \emph{built} once arbitrage is
removed.
Its existence \emph{is} the removal.
A single measure of what a computation has is not an additional fact
about the world; it is the shadow cast by a conversion structure with no
cycles that pay.
\end{remark}

\begin{remark}[Gauge]
\label{rmk:vtok-gauge}
The freedom in $p$ up to an additive constant is the freedom to choose
where zero sits.
Naming an existing token type as the reference --- pricing in one of the
colours, as one prices in dollars --- fixes the gauge at a vertex.
The virtual token is the refusal to privilege any vertex.
\end{remark}

\section{When arbitrage is present}

When $\hypK$ is not arbitrage-free, $\ell$ is not exact, and the graph
Hodge decomposition splits it into an exact part and a cyclic part:
\[
  \ell \;=\; \underbrace{-\,\delta p}_{\text{the valuation that does exist}}
       \;+\; \underbrace{\ell^{\mathrm{cyc}}}_{\text{the arbitrage}} .
\]
The first summand is the best common measure the system admits; the
second is a cohomology class, and it measures precisely how much the
system fails to have one.
It is not noise.
It is a quantity, it is computable from the rate data, and a computation
that can find a cycle with yield greater than one can profit from it
without doing any work at all.

Part~\ref{part:physics} identifies the exact part with energy and the
cyclic part with the failure of energy to be a state function, and
Chapter~\ref{sec:cost} draws the consequences: a first law, a second law,
and a bound on erasure.
Chapter~\ref{ch:eng} does something different with the same
decomposition --- it observes that when several communities of learners
each maintain their own token and must convert between them, the cycles
of the conversion graph carry redundancy, and redundancy is an
error-correcting code.
The reader who wants the physics should go to the former and the reader
who wants the sociology to the latter.
Both are reading the same short theorem.

% ===================================================================
